\section{Evaluation protocol}
\label{sec:protocol}

Backtested strategies are routinely overfitted. When many variants are tried on the same data,
the best one looks good partly by luck \citep{harvey2016,bailey2014}. We therefore fixed the
following protocol before running any experiment.

\begin{enumerate}
  \item \textbf{Split.} Data before \SplitDate{} are in-sample and used for all design and
        parameter choices. Data from \SplitDate{} to \DataEnd{} are the held-out test.
  \item \textbf{Enforcement by tooling.} The experiment runner reports in-sample statistics only.
        Test-period statistics appear only when the flag \code{--reveal-test} is passed, which will
        happen once, for a small set of final variants chosen in advance. The report-generation
        script cuts the data at the split date before it does anything else.
  \item \textbf{Trial counting.} Every variant run is saved with its configuration. At the time of
        writing, \NVariants{} distinct variants have been evaluated in-sample
        (\cref{app:variants}). This count will enter the deflated Sharpe ratio of the final
        results.
  \item \textbf{Fixed costs.} All variants are compared at the same one-way cost of 5~bps, the
        value used by AL. The break-even cost $c^\ast$ (\cref{eq:breakeven}) shows how sensitive
        each conclusion is to that assumption.
\end{enumerate}

\todo{Add the deflated Sharpe ratio, block-bootstrap confidence intervals for the Sharpe ratio,
and cost-sensitivity curves.}
